
% !TEX program = lualatex
\documentclass[a4paper,12pt]{bxjsbook}

% =========================================================
% Packages
% =========================================================

% --- Japanese (LuaLaTeX) ---
\usepackage{luatexja}

% --- math / layout ---
\usepackage{amssymb}
\usepackage{amsthm}
\usepackage{mathtools}
\usepackage{geometry}
\usepackage{mathpartir}
\geometry{margin=25mm}

% --- lists / links / figures ---
\usepackage{enumitem}
\usepackage{xurl}
\usepackage{hyperref}
\usepackage{tikz}
\usetikzlibrary{arrows.meta, positioning, calc, fit, backgrounds, decorations.markings}

\hypersetup{
  colorlinks=true,
  linkcolor=blue,
  urlcolor=blue,
  citecolor=blue
}

% =========================================================
% length definition
% =========================================================
\newlength{\customindent}
\setlength{\customindent}{2em}

% =========================================================
% macro definition
% =========================================================
\newcommand{\BodyListFixBegin}{%
  \begingroup
  \setlength{\rightskip}{0pt}%
  \setlength{\parindent}{0pt}%
  % \setlength{\displaywidth}{\linewidth}%
}
\newcommand{\BodyListFixEnd}{\ifhmode\par\fi\endgroup}

\newenvironment{BodyListFix}{\BodyListFixBegin}{\BodyListFixEnd}

% =========================================================
% environment definition
% =========================================================

% Main_Sentense styles
\newenvironment{body}{
  \par
  \begingroup
  \setlength{\leftskip}{\customindent}
  \setlength{\rightskip}{0pt}
  \setlength{\parindent}{1em}
  \setlength{\displayindent}{\customindent}
  \setlength{\displaywidth}{\dimexpr\linewidth-\customindent\relax}
}{
  \par
  \endgroup
}

% description in Main_Sentense styles
\newlist{bodydescription}{description}{1}
\setlist[bodydescription]{
  labelwidth=0pt,
  labelsep=0pt,
  itemindent=0pt,
  style=nextline,
  font=\bfseries
}

% itemize in Main_Sentense styles
\newlist{bodyitemize}{itemize}{3}
\setlist[bodyitemize]{
  leftmargin=3em,
  label=\textbullet,
  before=\begin{BodyListFix},
  after=\end{BodyListFix}
}
\setlist[bodyitemize,2]{leftmargin=3.5em}
\setlist[bodyitemize,3]{leftmargin=4em}

% enumerate in Main_Sentense styles
\newlist{bodyenumerate}{enumerate}{3}
\setlist[bodyenumerate]{
  leftmargin=4em,
  label=(\arabic*),
  before=\begin{BodyListFix},
  after=\end{BodyListFix}
}
\setlist[bodyenumerate,2]{leftmargin=4.5em, label=(\alph*)}
\setlist[bodyenumerate,3]{leftmargin=5em, label=(\roman*)}

% Theorem styles
\newtheoremstyle{theorem}
  {\baselineskip}
  {\baselineskip}
  {
    \normalfont
    \setlength{\leftskip}{\customindent}
    \setlength{\parindent}{0pt}
    \setlength{\rightskip}{0pt}
    \setlength{\displayindent}{\customindent}
    \setlength{\displaywidth}{\dimexpr\linewidth-\customindent\relax}
  }
  {0pt}
  {\bfseries}
  {.}
  {\newline}
  {
    \hskip-\customindent
    \thmname{#1}\thmnumber{ #2}\thmnote{（#3）}
  }

\theoremstyle{theorem}
\newtheorem{theorem}{定理}[chapter]
\newtheorem{lemma}[theorem]{補題}
\newtheorem{proposition}[theorem]{命題}
\newtheorem{corollary}[theorem]{系}
\newtheorem{definition}[theorem]{定義}
\newtheorem{example}[theorem]{例}
\newtheorem{remark}[theorem]{注意}
\newtheorem{abbreviation}[theorem]{略記}

% description in Theorem styles
\newlist{theoremdescription}{description}{3}
\setlist[theoremdescription]{
  leftmargin=5em,
  label=\textbullet,
  before=\begin{BodyListFix},
  after=\end{BodyListFix}
}
\setlist[theoremdescription,2]{leftmargin=5.5em}
\setlist[theoremdescription,3]{leftmargin=6em}

% itemize in Theorem styles
\newlist{theoremitemize}{itemize}{3}
\setlist[theoremitemize]{
  leftmargin=5em,
  label=\textbullet,
  before=\begin{BodyListFix},
  after=\end{BodyListFix}
}
\setlist[theoremitemize,2]{leftmargin=5.5em}
\setlist[theoremitemize,3]{leftmargin=6em}

% enumerate in Theorem styles
\newlist{theoremenumerate}{enumerate}{3}
\setlist[theoremenumerate]{
  leftmargin=6em,
  label=(\arabic*),
  before=\begin{BodyListFix},
  after=\end{BodyListFix}
}
\setlist[theoremenumerate,2]{leftmargin=6.5em, label=(\alph*)}
\setlist[theoremenumerate,3]{leftmargin=7em, label=(\roman*)}

% =========================================================
% Proof Environment Setting
% =========================================================
\makeatletter
\renewenvironment{proof}[1][\proofname]{%
  \par\pushQED{\qed}%
  \normalfont
  \topsep6\p@\@plus6\p@\relax
  \list{}{%
    \leftmargin=\customindent
    \rightmargin=0pt
    \labelwidth=0pt
    \labelsep=0pt
    \itemindent=0pt
    \listparindent=0pt
    \parsep=0pt
  }%
  \item[\itshape #1\@addpunct{.}\hspace{0.5em}]%
  \setlength{\displayindent}{\customindent}%
  \setlength{\displaywidth}{\dimexpr\linewidth-\customindent\relax}%
  \ignorespaces
}{%
  \popQED\endlist\@endpefalse
}
\makeatother

% =========================================================
% Convenience commands
% =========================================================
\newcommand{\addchaptertotoc}[1]{%
  \chapter*{#1}%
  \addcontentsline{toc}{chapter}{#1}%
}
\newcommand{\dotminus}{\mathbin{\dot{-}}}

\DeclareMathOperator{\AffHull}{AffHull}
\DeclareMathOperator{\Span}{Span}
\DeclareMathOperator{\AffSpan}{AffSpan}
\DeclareMathOperator{\AffComb}{AffComb}
\DeclareMathOperator{\ConvHull}{ConvHull}
\DeclareMathOperator{\Face}{Face}
\DeclareMathOperator{\Facet}{Facet}


% =========================================================
% Metadata
% =========================================================
\title{チェバの定理の一般化と\\ Lean4/mathlib4への形式化}
\author{
秋田 隼\\
早稲田大学\\
基幹理工学部 数学科B4\\
1W221002-0
}

\date{$2026$/$1$/$9$}

% =========================================================
% Document
% =========================================================
\begin{document}

    \maketitle

    \addchaptertotoc{進捗(最終的に消します)}
    \textbf{必要十分条件を主張する定理として一般化を行いたかったが、どのバージョンの必要十分形が最も汎用的で理想のチェバの定理と言えるのか結論が出ず、また仮に結論が出たとして、おそらく実装が間に合わないことから、ジョセフマイヤーズ氏の定理をそのまま逆にした定理を示し実装することにしました。もし余裕があればさらに一般化を行いたいと思います。}
    2026/1/22 2.3 Lean4・mathlib4の説明を編集
    2026/1/30 一般化とは何か、古典的チェバの主張になぜパターンがあるかの説明を補足、定義の整合性や証明部分はまだ未完成です

    \frontmatter
    \addchaptertotoc{要旨}

    \begin{body}
        \indent 本研究は、古典的チェバの定理を出発点として、チェバの定理の一般化を定式化し、関連する補題のLean4/mathlib4上での形式化を目標とする。

        \vskip\baselineskip

        \indent 主な貢献は以下である。

        \begin{bodyitemize}[leftmargin=3em]
            \item チェバの定理の一般化の定式化の紹介。
            \item 既存のmathlib4の構造を調査し、再利用可能部分と不足部分を切り分ける設計指針。
            \item 不足する補題・定義をモジュール化して追加し、定理の機械検証を通す。
        \end{bodyitemize}

        \vskip\baselineskip
        \indent 実装はGitHubリポジトリで公開している：
        \par\noindent
        \url{https://github.com/Aj1905/LeanResearch.git}
    \end{body}

    \tableofcontents

    \mainmatter

    % =========================================================
    \chapter{序論}
    % =========================================================

    \section{研究背景と動機}
    \begin{body}
        \indent 本研究は、mathlib4に存在する順方向チェバの定理の形式化を基盤として、その定理の逆をLean4上で構成し、mathlib4への追加を目指す。mathlib4は、Lean Prover Communityにより開発される大規模数学ライブラリ\cite{LeanProverCommunity}であり、分野により整備状況には濃淡がある。
        特に幾何学は、形式化が代数的手法に比べて技術的に困難であることや、図形的直観の形式言語翻訳の複雑さが起因するためか発展が停滞気味である。こうした開発目標を可視化すべく、数学において重要でありながら未だ形式化されていない定理が「missing 100 theorems」としてリストアップされている。

        \indent チェバの定理は、図形の共点性を代数条件へ翻訳する基本的な道具であり、多数の幾何学的議論の分岐点として機能する。この種の定理を一般化してライブラリ化することの価値は、定理本体に加えて単体・面・アフィン結合・重心座標などの重要な周辺概念を再利用可能な形で整備できる点にある。よって今後の幾何定理を形式化する際の基盤を構築できるという点で大きな意義を持つと考える。

        \indent さらに、形式化は教科書的な証明に潜む、ゼロ除算回避、図形が退化する場合の排除、必要な代数構造などの暗黙の仮定を明示し、どの仮定のもとでどの結論が成立するかを定理の型として確定する。このことは、一般化の汎用性を証明の文章ではなくLeanの型と依存関係によって機械的に検証可能な形で提示できる点で、数学的にも工学的にも意義がある。近年もチェバ型定理の一般化は研究されており、本研究はその形式的基盤をmathlib4に提供する試みとして位置づけられる。
    \end{body}

    \section{本稿の構成}
    \begin{body}
        \indent 第$2$章でLean4について解説し、第$3$章で一般化とは何か基礎論的な視点から眺める。第$4$章でチェバの定理とアフィン幾何の関連を議論し、第$5$章で古典的チェバの定理およびその逆がどのように一般化され得るかを考える。第$6$章で実際に行った実装の詳細を見る。
        \indent 本稿の議論は暗黙に、自然演繹を推論規則として採用した古典一階述語論理体系を用いて行なっている。集合論に言及する時、その概念の定義はKenenに従うものとする。\cite{Kunen}
    \end{body}


    % =========================================================
    \chapter{Lean4による形式化}
    % =========================================================
    \section{形式言語の概要}
    \begin{body}
        \indent 自然言語の主張はさまざまな点で曖昧さを含む。定義が省略されていたり、暗黙の前提があったり、同じ記号が別の意味を持っていたりすることで本来したいはずの主張が伝わらないことや別の意味で解釈されてしまうこともある。

        \indent これを解決するために用いられるのが\emph{形式言語}と呼ばれる定義が厳密に定まった言語体系であり、自然言語の文を形式言語に翻訳する作業は\emph{形式化}と呼ばれる。一度形式化された文章は、読み手がその体系の定義を正確に把握している限り、いつ誰がどこで読んでも意味が一意に定まる。これはコンピュータが文章を読む場合も例外ではない。この形式言語の文章をどう解釈しどう推論するかも含めて厳密に定義したものを形式体系と呼ぶ。

        \indent 一方数学とは、過去の数学者の成果の積み重ねからなる学問である。数学的主張は、他者に共有され、検証され、再利用される形で後世に伝え残されてはじめて真価を発揮する。このとき、時代や場所、読み手によって意味や解釈が揺れることは望ましくない。いついかなる時も同じ主張をすることが求められる。まさに形式言語の特徴が最大限活かされる分野の一つである。現代数学は集合論を一階述語論理で公理化した形式体系（主にZFやZFC）を基礎として用いて構築されことが多い。

        \indent しかし一度考え直してみると、数学を記述する形式言語は必ずしも一階述語論理に限る必要はない。一階述語論理は確かに人間が数学を理解する上で相性のいい形式言語の一種かもしれない。形式言語の中では比較的読みやすく、また歴史的に標準化されているため慣れてもいる。しかし本来、言語は読み手と目的に最適化することで意味伝達という言語本来の役割を最大限発揮する。伝えたい意味内容に対し、読み手と目的を明確にして最適な形式言語を選択することが重要となる。そのための一階述語論理の代替案の一つとして挙げられるのが型理論である。

        \indent 型理論とは、命題を型、証明をその型の項として表現するという特徴を持った形式体系である。集合論において全ての数学的対象が集合であるように、型理論において全ての数学的対象は型を持つ。詳細な説明は\hyperref[appendix:type-theory]{付録}にまとめた。たとえば命題 $P$ の証明は $p:P$（pは型Pの項）という形で与えられる。項$p$を構築する厳密な規則が用意されているため、その規則に従っているかを機械的に確認する作業を通して証明の正しさを検証できる。その意味で、コンピュータにとって検証のしやすい言語であると言える。

    \end{body}

    \section{証明支援系の概要}
    \begin{body}
        \indent 証明支援系とは、数学の定義・定理・証明を形式言語で記述することに特化したプログラミング言語、かつその言語で記述されたコードが公理や推論規則に照らして正しいかを機械的に検証する機能を備えたもの、の総称である。
        
        \indent 「証明を自動で発見する」ことを目的としたATP（Automated Theorem Proving）と、「与えられた証明が正しいことを厳密に検査する」ことを目的としたITP（Interactive Theorem Proving）が存在する。

        \indent 証明支援系には基礎となっている形式体系の種類に応じて複数の系統がある。代表的なものとして、依存型理論に基づくLeanやCoq（現在Rocqに改名）、高階論理系のIsabelle、HOLなどが存在し、目的に応じて使い分けられている。

        \indent この技術は応用面では、まずソフトウェアの形式検証に用いられる。ソフトウェアは各々が仕様という設計目標を持つ。この仕様を数学的な命題として書き下すことで、実装されたプログラムがその仕様を満たしているか否かという問題を、数学的な真偽判定の問題に帰着させることができる。挙動をあらかじめ検証しておくことで、信頼性の高い堅牢なソフトウェア開発が可能となる。

        \indent また近年、証明支援系はLLMとも結びつき始めている。LLMはIMO 2025(国際数学オリンピック)において金メダル基準相当の成績に達したとする報告も出ており\cite{DeepMindIMO2025}、自然言語推論の性能が飛躍的に向上している。一方で、もっともらしいが誤りを含む出力（ハルシネーション）は依然としてLLMの課題である。
        しかし、近年注目されている定理証明器を統合した枠組みでは、自然言語の解答案全体または一部を形式化して検証し、失敗時にはフィードバックにより解答を再生成するというループが可能となる。そのため、形式検証が及ぶ範囲についてはハルシネーションを大幅に抑制でき、結果として出力全体に対するハルシネーションの頻度も抑えられる。これは、ハルシネーション自体を抑えるのではなく正確な解答を返すまで出力を繰り返させるという発想の転換である
        （下図参照、問題文の形式化や自然言語への説明生成には依然として誤りが入りうる）。よって、「形式検証の及ぶ範囲を拡大する」、「問題文の形式化精度を向上させる」ことで生成AIの数学力は向上すると考えられており、研究が盛んに行われている。

    \end{body}

    \begin{tikzpicture}
        %========================
        % 共通
        %========================
        \pgfmathsetlengthmacro{\Span}{0.90*\linewidth}

        \pgfmathsetlengthmacro{\TopY}{0pt}
        \pgfmathsetlengthmacro{\BotY}{-3.5cm}

        \coordinate (TopL) at (-0.5*\Span,\TopY);
        \coordinate (TopR) at ( 0.5*\Span,\TopY);
        \coordinate (BotL) at (-0.5*\Span,\BotY);
        \coordinate (BotR) at ( 0.5*\Span,\BotY);

        %========================
        % 上段
        %========================
        \tikzset{
        box/.style={
        rectangle, draw, thick,
        text width=3.0cm,
        minimum height=1.0cm,
        align=center, font=\small,
        inner sep=2pt
        },
        myarrow/.style={->, >=Stealth, thick, shorten <=2pt, shorten >=2pt},
        danger/.style={box, fill=red!20},
        }

        \node[box] (nlProblem)   at ($(TopL)!0.00!(TopR)$) {自然言語\\問題};
        \node[box] (nlReasoning) at ($(TopL)!0.42!(TopR)$) {自然言語\\推論};
        \node[box] (nlAnswer)    at ($(TopL)!0.84!(TopR)$) {自然言語\\解答};

        \draw[myarrow] (nlProblem) -- (nlReasoning);
        \draw[myarrow] (nlReasoning) -- (nlAnswer);

        \node[above=3mm of nlProblem] {\bfseries 従来の自然言語推論};

        %========================
        % 下段
        %========================
        \tikzset{
        boxS/.style={
        rectangle, draw, thick,
        text width=1.85cm,
        minimum height=0.72cm,
        align=center,
        font=\scriptsize,
        inner sep=2pt
        },
        dangerS/.style={boxS, fill=red!20},
        safeS/.style={boxS, fill=green!20},
        myarrowS/.style={->, >=Stealth, thick, shorten <=2pt, shorten >=2pt}
        }

        \node[boxS]    (flProblemNl) at ($(BotL)!0.00!(BotR)$) {自然言語\\問題};
        \node[safeS]   (flProblem)   at ($(BotL)!0.21!(BotR)$) {形式言語\\問題};
        \node[safeS]   (flReasoning) at ($(BotL)!0.42!(BotR)$) {形式言語\\推論};
        \node[safeS]   (flAnswer)    at ($(BotL)!0.63!(BotR)$) {形式言語\\解答};
        \node[boxS] (finalAnswer) at ($(BotL)!0.84!(BotR)$) {自然言語\\解答};

        \draw[myarrowS] (flProblemNl) -- (flProblem);
        \draw[myarrowS] (flProblem)   -- (flReasoning);
        \draw[myarrowS] (flReasoning) -- (flAnswer);
        \draw[myarrowS] (flAnswer)    -- (finalAnswer);

        \begin{pgfonlayer}{background}
            \node[draw, thick, inner sep=2pt, rounded corners=2pt, fit=(flProblem)(flReasoning)(flAnswer), green, label={[font=\scriptsize]above:ハルシネーションを抑制}] (flBigBox) {};
        \end{pgfonlayer}

        \draw[myarrowS, dashed]
        (flReasoning.south) .. controls +(0,-0.8) and +(0,-0.8) ..
        node[pos=0.55, right, font=\tiny, xshift=-2mm, yshift=-2mm] {検証失敗}
        (flProblem.south);

        \node[above=3mm of flProblemNl] {\bfseries 形式検証統合推論};

    \end{tikzpicture}

    \begin{body}
        \indent 上記を踏まえると本稿の目的は、\emph{数学という体系（特にチェバの定理）を証明支援系を通してコンピュータに検証してもらうために適切な形式言語を選択しコンピュータに伝達すること}である。

        \indent この手段として私はLean4という言語を選択した。
    \end{body}

    \section{Lean4・mathlib4の概要}
    \begin{body}
        \indent Lean4とは、依存型理論の一種CIC（Constructive Inductive Calculus）を基礎理論とする、プログラミング言語であり証明支援系（ITP）でもある。

        \indent この言語の概要を知るにはまずLean4のファイル構成を把握するのが良いが、これは一般的な数学書と大きく変わらない。論理の流れに従い$def$、$lemma$や$theorem$という宣言に続いて用語の定義、命題の主張が並んでいる。数学書が一部の定理の証明を他の文献に譲るように、Lean4のファイルでは別のファイルを参照し、必要な定理をインポートすることでその定理を既知の定理として証明中で再利用できるようになる。再利用前提でファイルを作成する必要があるため、ファイル構成にもある程度の形式・ルールが生まれる。この要件を満たしながら定義・定理の数々をLean4に書き起こして行ったファイルを集めたものが、前述したmathlib4という巨大なライブラリである。証明支援系、ライブラリという人によっては馴染みのない用語を聞いて身構える読者もいるかもしれないが、Lean4を一種のプログラミング言語として捉えてしまえば、数学体系をLean4という言語で書き直し、mathlib4という共用のフォルダにまとめて管理しているだけである。世界中の数学者や技術者が共同でこのフォルダを管理・運営し、世界規模で開発が行われている。この言語は2021年にLean3を改良して生まれた新しい言語であるが、両者に互換性がないため、現在Lean3時代のライブラリ（mathlib3）からのLean4への翻訳が進められている。

        \indent そんなLean4であるが、数学書と最も大きく異なる点がその言語が型理論をベースにして構築されていることである。一般的な数学書は一階述語論理をベースに部分的に自然言語による説明を交えた形式で綴られるが、Lean4ではそのすべてが一つの言語を用いて記述される。数学的主張を世に出す立場として、慣れるまで少し時間がかかるのはデメリットであるが、一方でLean4が持つ大きなメリットが「証明を書く過程でその証明が正しいかどうかを判定してくれる点」である。背後で依存型理論をベースとした型チェッカーというツールが稼働し続け、記述している証明が誤っていれば直ちにエラーを返してくれる場合があり、非常に効率的に記述が行える。いわば、"添削機能付きの数学書"である。

        \indent Lean4における証明記述には大別して2つのモードが存在する。
        
        \begin{bodyitemize} [leftmargin=1.8em]
            \item[]
                \textbf{項モード}\par\noindent
                証明に相当する項をコーディングしながら自分で構成していく。
            \item[]
                \textbf{tacticモード}\par\noindent
                あらかじめ定義されたtacticと呼ばれるコマンドを利用して、証明すべき命題を変更するなどの操作を行う。
        \end{bodyitemize}

        \vskip\baselineskip

        \indent 詳細な説明は割愛するが、前者は証明全体の流れを明示しやすく，後者は手計算の証明手順に近い形で推論を整理できるため，両者は相補的に用いられる。
    \end{body}

    % =========================================================
    \chapter{構文論における定理の一般化}
    % =========================================================
    \begin{body}
        \indent この証明では一般化とはそもそも何か、定理の逆とは何かを構文論の観点から振り返る。
    \end{body}

    \section{定理の一般化とは}

    \begin{body}
        \indent 命題の集合$T$に対し、
        \[
        T \vdash \phi
        \]
        の形で表される命題$\phi$のことを（$T$から導出される）\emph{定理}と呼ぶ。特に$\phi$が$p\to q$の形のとき、命題$q\to p$を\emph{定理の逆}と呼ぶ。定理の逆は一般には真になるとは限らない。   

        \indent また、複数の命題をまとめて定理と呼ぶ場合も多い。$\phi_{1},\phi_{2},\cdots,\phi_{n}$が全て$T$から導出される定理である時、
        \[       
        T \vdash (\phi_{1}\land\phi_{2}\land\cdots\land\phi_{n})
        \]
        一方、自然演繹を推論規則として採用する上では演繹定理と呼ばれる以下の性質が成り立つ。
        \[
        T\cup\{\psi\}\vdash \phi \;\Longleftrightarrow\; T\vdash (\psi\to\phi).
        \]
        $T= \{\varphi_{1},\varphi_{2},\cdots,\varphi_{n}\}\cup\{\psi_{1},\psi_{2},\ldots,\psi_{n}\}$の時、
        \[
        \begin{alignedat}{2}
        &\phantom{\iff}\; T \vdash \phi && \\
        &\iff \{\varphi_{1},\varphi_{2},\cdots,\varphi_{n}\}\vdash\bigl(\psi_{1}\to(\psi_{2}\to(\cdots\to(\psi_{n}\to\phi)))\bigr) &&\text{\footnotesize（演繹定理を繰り返し適用）}\\
        &\iff \{\varphi_{1},\varphi_{2},\cdots,\varphi_{n}\}\vdash\bigl(\neg\psi_{1}\lor(\neg\psi_{2}\lor(\cdots\lor(\neg\psi_{n}\lor\phi)))\bigr) &&\text{\footnotesize（含意と論理和の関係）}\\
        &\iff \{\varphi_{1},\varphi_{2},\cdots,\varphi_{n}\}\vdash\bigl(\neg\psi_{1}\lor\neg\psi_{2}\lor\cdots\lor\neg\psi_{n}\lor\phi\bigr) &&\text{\footnotesize（$\lor$の結合法則）}\\
        &\iff \{\varphi_{1},\varphi_{2},\cdots,\varphi_{n}\}\vdash\bigl(\neg(\psi_{1}\land\psi_{2}\land\cdots\land\psi_{n})\lor\phi\bigr) &&\text{\footnotesize（ド・モルガンの法則）}\\
        &\iff \{\varphi_{1},\varphi_{2},\cdots,\varphi_{n}\}\vdash\bigl((\psi_{1}\land\psi_{2}\land\cdots\land\psi_{n})\to\phi\bigr) &&\text{\footnotesize（含意と論理和の関係）}\\
        \end{alignedat}
        \]
        よって、$\phi_{1},\phi_{2},\cdots,\phi_{n}$が$T$から導出される定理である時、以下のように書き下せる。
        \[
        \{\varphi_{1},\varphi_{2},\cdots,\varphi_{n}\}\vdash{\psi_{1}\land\psi_{2}\land\cdots\land\psi_{n}}\to{\phi_{1}\land\phi_{2}\land\cdots\land\phi_{n}}
        \]
        
        \vskip\baselineskip

        $P:=\{\varphi_{1},\varphi_{2},\cdots,\varphi_{n}\},A:=\{\psi_{1},\psi_{2},\cdots,\psi_{n}\},C:=\{\phi_{1},\phi_{2},\cdots,\phi_{n}\}$と置き、便宜上以下のように略記する
        \[
        P\vdash A\to C
        \]
        これは自然演繹の公式の記法ではなく、説明上の都合でこのような書き方になっていることに留意してほしい。
        この$P$を\emph{前提/背景理論}、$A$を\emph{仮定}、$C$を\emph{結論}と呼ぶ。
        要するに、定理が与えられるとそれを適当な粒度の命題の組を用意し、\emph{前提・仮定・結論の導出関係として書き換えることができる}ということである。

        \indent 数学基礎論以外の分野の一般的な数学書では前提$P$として$ZF/ZFC$が慣習的に用いられるため、これらは暗黙に受け入れたものとして敢えては明記されず、仮定$A$から結論$C$を導く、というステップのみが記述されることが多い。

        \indent また、この形式で定理を表現することにより、暗黙に認めていた前提を明文化できるため、Lean4での形式化に繋げやすいというメリットがある。

        \indent この上で、本稿で言う「定理の一般化」とは、既存の定理がある意味で\emph{特別な場合として含まれる}ような、より広い主張をする定理を与えることである。典型的には次の三方向がある。
        \begin{bodyitemize}[leftmargin=3.5em]
            \item \textbf{前提の弱化}：$P$の元の命題の主張を弱める。$\varphi\in P$を、$\varphi\to \varphi'$を満たす$\varphi'$で置き換える。 
            \item \textbf{仮定の弱化}：$A$の元の命題の主張を弱める。$\psi\in A$を、$\psi\to \psi'$を満たす$\psi'$で置き換える。
            \item \textbf{結論の強化}：$C$の元の命題の主張を強める。$\phi\in C$を、$\phi'\to \phi$を満たす$\phi'$で置き換える。
            \item \textbf{複数の組合せ}：上記を複数個組み合わせる。
        \end{bodyitemize}
        
        \indent (3)は「定理の強化」と呼べれることが多いが、本稿では一般化に含める。一般化と聞くと仮定と結論に着目しがちであるが、定理が依存する前提を緩め（例：体から環へ、可換環から一般の環へ、有限次元から任意次元へ）、対象の主張をより広い土台で成立させることも一般化の重要な形である。

    \end{body}

    % =========================================================
    \chapter{チェバの定理とアフィン空間}
    % =========================================================
    \begin{body}
        \indent チェバの定理とは高校数学でも頻繁に登場する、初等幾何における重要な定理の一つである。

        \indent 本章ではこのチェバの定理の詳細と、チェバの定理が展開されるアフィン空間の性質について議論する。

    \end{body}

    \section{チェバの定理の有名な形}

    \begin{body}
        \indent 高校の教育課程において、チェバの定理は以下の形で紹介されることが多い。

        \indent ただし、定理中で登場する用語は初等幾何的な高校の学習課程内の意味である。
    \end{body}

    \begin{theorem}[古典的チェバの定理]\label{thm:classical-ceva}
        三角形$\triangle ABC$の各辺$BC,CA,AB$またはその延長線上にそれぞれ頂点とは異なる点$D,E,F$をとる。
        このとき$3$直線$AD$、$BE$、$CF$が$1$点で交わるならば、以下の式が成立する
        \[
          \dfrac{BD}{DC}\,\dfrac{CE}{EA}\,\dfrac{AF}{FB} = 1
        \]
    \end{theorem}

    \begin{theorem}[古典的チェバの定理の逆]
        三角形$\triangle ABC$の各辺$BC,CA,AB$またはその延長線上にそれぞれ頂点とは異なる点$D,E,F$をとる。
        このとき
        \[
          \dfrac{BD}{DC}\,\dfrac{CE}{EA}\,\dfrac{AF}{FB} = 1
        \]
        が成立するならば、$3$直線$AD,BE,CF$は$1$点で交わる、または互いに平行である。
    \end{theorem}

    \begin{body}
        \indent ここで登場する点$D,E,F$を\emph{チェビアンの足}、$3$直線$AD,BE,CF$を\emph{チェビアン}と呼ぶ。

        \indent つまり古典的チェバの定理（またはその逆）とは\emph{チェビアンが共点するという幾何的性質が成立する条件を代数的に記述する定理}である。

        \indent 上記のチェバの定理の主張には、一見すると線分の長さ（BD、DC等）が現れるため、ユークリッド幾何に依存する定理のように見える。

        \indent しかし、実際にはこの定理の主張は線分の長さの比（$\dfrac{BD}{DC}$ 等）のみで記述されている。実はこれらの概念は後述するようにユークリッド幾何を用いることなくアフィン幾何という体系内で定義することができる。

        \indent またチェバの定理が必要十分性を主張する定理だと認識していた読者も多いかもしれない。$D,E,F$を辺上に限定する版や、$3$直線が$1$点で交わる必要十分条件を提供する版など、教科書により紹介の仕方に微妙な違いがあり、これらは教育的配慮によるものである。詳細は\hyperref[appendix:classical-ceva-version]{付録}で解説するが、アフィン幾何の範疇では古典的チェバの定理の主張したい内容が「共点$⇔$等式条件」というシンプルな形の必要十分条件では表せないという事にだけ留意してほしい。

        \indent 以降、アフィン幾何とはどのような体系なのか紹介する。
    \end{body}

    \section{概念の定義}
    \begin{body}
        \indent まず今後の議論で登場する、数学的概念を定義する。定義の多くは紙面上で定理の証明を行いやすくなるものを採用した。

        \indent また、定義から導かれるいくつかの補題を$remark$欄に紹介している。本稿の主題ではないのでこれらは認めるものとして今後議論を進める。詳細な証明はmathlib4ライブラリや、アフィン幾何学の教科書を参考にしてほしい。

        \indent 以下特に断らない限り、$k$は環、$V$は$k$-加群、$\iota$は有限の添字集合とする。
    \end{body}

    \begin{definition}[アフィン空間]
        $P$ が $V$ 上の（$k$ に関する）\emph{アフィン空間}であるとは、次の2つの写像が存在し、以下の公理を満たすことをいう：
        \[
        \dotplus :\; V\times P \ni (v,p) \mapsto v \dotplus p \in P,
        \qquad
        \dotminus:\; P\times P \ni (p,q) \mapsto p \dotminus q \in V.
        \]
        \begin{theoremitemize}[leftmargin=5em, topsep=2pt, partopsep=0pt, itemsep=0pt, parsep=0pt]
            \item[\hspace*{5em}\textbf{公理1}] 任意の $p\in P$ と $v,w\in V$ について
            \[
            \mathbf{0}\dotplus p = p.
            \]
            \item[\hspace*{5em}\textbf{公理2}] 任意の $p,q\in P$ について
            \[
            (p\dotminus q) \dotplus q= p.
            \]
            \item[\hspace*{5em}\textbf{公理3}] 任意の $p\in P$ と $v\in V$ について
            \[
            (v\dotplus p)\dotminus p = v.
            \]
        \end{theoremitemize}
        アフィン空間 $P$ の元を\emph{点}と呼ぶ。
        写像$p:\iota \to P$が定義する点の集合$Im(p)$、またはその写像自体を\emph{点族}と呼ぶ。
        しばしば点族の部分集合を$s\subseteq \iota$を用い、$p$の制限として$p|_s$と表す。
    \end{definition}

    \begin{remark}[注意：点同士の和は一般に定義されない]
        上の定義で与えられるのは「点 $p\in P$ に加群の元 $v\in V$ を足す演算 $p\dotplus v$」と
        「2点 $p,q\in P$ の差から加群の元を作る演算 $q\dotminus p$」である。
        アフィン空間とは、非公式には「原点を忘れたベクトル空間」と説明されることが多い。要は、「点を平行移動させるとまた点である」、「2点の間には対応する平行移動がある」を満たす空間である。
        ベクトル空間の基本的な演算である「$p,q\in P$ の和 $p+q$」や「スカラー倍 $\lambda p$」に相当する演算は、アフィン空間中そのままでは定義されない。
    \end{remark}

    \begin{remark}[注意 : $V$に$k$-加群が必要な理由]
        純粋なアフィン空間を定義するだけなら$V$は加法群で足りる。
        公理のいずれにもスカラー倍の概念は現れない。
        しかし以降導入するアフィン結合や重心座標ではスカラー係数を用いるため、$V$が$k$-加群である必要がある。
    \end{remark}

    \begin{abbreviation}[加法・減法（アフィン空間）]
        以下、簡単のため $\dotplus$ は $+$、$\dotminus$ は $-$ と表す。
    \end{abbreviation}

    \begin{definition}[アフィン結合（\texorpdfstring{$\AffComb$}{AffComb}）]
        写像 $\lambda:s\to k$ が
        \[
          \sum_{i\in s}\lambda(i)=1
        \]
        を満たすとき，任意に $j\in s$ を固定して
        \[
          \AffComb\bigl(p|_s,\lambda\bigr)
          \;:=\;
          p(j)+\sum_{i\in s}\lambda(i)\,\bigl(p(i)-p(j)\bigr)
        \]
        と定める。
        この表現形式を（点族 $p|_s$ の）\emph{アフィン結合}と呼び，$\lambda$ を\emph{重み}と呼ぶ。
    \end{definition}
      
    \begin{remark}[注意：基点 $j$ に依らないことと略記]
        条件 $\sum_{i\in s}\lambda(i)=1$ のもとでは，上の定義は $j\in s$ の選び方に依らない。
        以降は形式的な略記として
        \[
          \AffComb\bigl(p|_s,\lambda\bigr)=\sum_{i\in s}\lambda(i)\,p(i)
        \]
        と書くことがある。ただし右辺は「点の線形結合」という本来未定義である演算を用いた略記になっていることに留意してほしい。
    \end{remark}

    \begin{definition}[アフィン包]
        \[
          \AffHull(p,s)
          :=
          \left\{\,
            \AffComb\bigl(p|_s,\lambda\bigr)
            \ \middle|\
            \lambda:s\to k,\ \sum_{i\in s}\lambda(i)=1
          \,\right\}
        \]
        で定まる集合を点族 $p|_s$ の \emph{アフィン包}と呼ぶ。
        特に $s=\iota$（$\iota$ が有限である場合）には $\AffHull(p):=\AffHull(p,\iota)$ と略記する。
    \end{definition}

    \begin{definition}[アフィン部分空間]
        部分集合 $S\subseteq P$ が（$P$ の）\emph{アフィン部分空間}であるとは，
        $S\neq\emptyset$ であり，任意の有限集合 $t\neq \emptyset$、点族$q:t\to S$、重み $\mu:t\to k$ が
        \[
        \sum_{j\in t}\mu(j)=1 \Longrightarrow \AffComb(q,\mu)\in S
        \]
        が成り立つことをいう。
    \end{definition}
      
    \begin{remark}[注意：定義がスカラー係数を含む理由]
        アフィン部分空間は本来，$V$ の加法群作用（平行移動）だけで定義でき、スカラー倍を仮定する必要はない。
        しかし本稿ではアフィン結合・重心座標を中心に議論を進めるため、スカラー倍を前提にした定義になっていることに留意してほしい。
    \end{remark}

    \begin{definition}[直線]
        アフィン部分空間 $S\subseteq P$ が \emph{直線}であるとは，
        添字集合$t=\{0,1\}$, 点族$q:\iota\to S$に対して
        \[
        S = \AffHull(q,t)
        \]
        が成り立つことを言う。
    \end{definition}

    \begin{definition}[平行]
        アフィン部分空間 $S,T\subseteq P$ に対し、$T$が$S$に \emph{平行}であるとは，
        ある $v\in V$ が存在して
        \[
        T=\{\,v+p \mid p\in S\,\}
        \]
        と書けることをいう。このとき $S\parallel T$ と書く。
    \end{definition}

    \begin{definition}[アフィン独立]
        $s\subseteq \iota$ は空でないとする。
        点族 $p|_s:s\to P$ が \emph{アフィン独立}であるとは，
        任意の $i\in s$に対し、
        \[
        \sum_{j\in \iota\setminus\{i\}} c_j\,\bigl(p(j)-p(i)\bigr)=0
        \qquad\Longrightarrow\qquad
        \forall j\in \iota\setminus\{i\},\ c_j=0
        \]
        が成り立つことをいう。
    \end{definition}

    \begin{remark}[最小性]
        $\AffHull(p,s)$ は $Im(p|_s)$ を含む最小のアフィン部分空間である。\\
        この結果は認めて利用する。
    \end{remark}

    \begin{definition}[重心座標]
        $p|_s$ がアフィン独立であるとする。
        任意の点 $q\in \AffHull(p,s)$ に対して，ただ一つの重み $\lambda:s\to k$ が存在して
        \[
          q=\AffComb\bigl(p|_s,\lambda\bigr),
        \]
        を満たす。この $\lambda$ を $q$ の（$p|_s$ に関する）\emph{重心座標}と呼ぶ。
    \end{definition}

    \begin{remark}[注意：凸性を扱うには $k$ の順序が必要]
        以下の「凸包」「単体（凸単体）」では不等式 $\lambda(i)\ge 0$ を用いるため、
        この段以降では $k$ が順序環であることを仮定する。
    \end{remark}

    \begin{definition}[凸包]
        \[
          \ConvHull(p,s):=
          \left\{\,
            \AffComb\bigl(p|_s,\lambda\bigr)
            \ \middle|\
            \lambda:s\to k,\ \forall i\in s,\ \lambda(i)\ge 0,\ \sum_{i\in s}\lambda(i)=1
          \,\right\}
        \]
        を 点族$p|_s$ の \emph{凸包}と呼ぶ。
        特に $s=\iota$ のとき $\ConvHull(p):=\ConvHull(p,\iota)$ と略記する。
        これは $\AffHull(p)$ の要素の中で非負係数のアフィン結合で表せる点全体である。
    \end{definition}

    \begin{definition}[単体（凸単体）]
        $p|_s$ がアフィン独立であるとき，
        $p|_s$ の凸包を \emph{$(|s|-1)$-単体}と呼ぶ。
    \end{definition}

    \begin{example}
        $s$ が $n+1$ 点集合のとき、アフィン独立な$p|_s$が定める凸包を $n$-単体と呼ぶ。
        特に異なる3点 $A,B,C$ が定める凸包は $2$-単体であり、\emph{三角形 $ABC$} と呼び、$\triangle ABC$ と表す。
    \end{example}

    \begin{remark}
        mathlib中で単体は、「アフィン独立な点族」として定義されており、本稿の意味での単体は順序を定義した後構成するものとなっている。
    \end{remark}

    \begin{definition}[面]
        $S$を$p|_s$が定める$n$-単体とする。\\
        $s_0\subseteq s$に対し、\\
        点族$p|_{s_0}$のアフィン包を（$S$の$s_0$が定める）\emph{面}\\
        点族$p|_{s_0}$の凸包を（$S$の$s_0$が定める）\emph{狭義の面}\\
        と呼ぶ。
        特に $|s_0|=k+1$ のとき、$p|_{s_0}$が定める面を（Sの）\emph{$k$-面}と呼ぶ。
    \end{definition}

    \begin{definition}[対向面（ファセット）]
        $S$を$p|_s$が定める$n$-単体とする。\\
        $S$の$s\setminus \{i\}$が定める面を、頂点 $p(i)$ 対向面（またはファセット）と呼ぶ。
      \end{definition}      

    \begin{example}
        $\triangle ABC$（$2$-単体）のとき、頂点 $A$ の対向面は辺 $BC$ である。
    \end{example}

    \begin{definition}[共点する]
        直線$S_1,\cdots,S_m$ が一点 $x\in P$ を共有するとき、
        すなわち
        \[
          x \in \bigcap_{i\in s} S_i
        \]
        を満たすとき、これらは \emph{共点する}という。
        また$x$を\emph{共点}と呼ぶ。
    \end{definition}

    % \begin{definition}[チェビアン / チェビアンの足(古典版)]
    %     $S$ を $2$-単体とし、$i\in s$ とする。
    %     頂点 $p(i)$ を通る直線 $S$ が
    %     \[
    %     L\cap \AffHull(p,s\setminus\{i\}) \neq \varnothing
    %     \]
    %     を満たすとき、この直線 $S$ を（頂点 $p(i)$ に関する）\emph{チェビアン}と呼ぶ。
    %     また、交点
    %     \[
    %     x:=L\cap \AffHull(p,s\setminus\{i\})
    %     \]
    %     が一点で定まるとき、その点 $x$ を \emph{チェビアンの足}と呼ぶ。
    % \end{definition}

    % \begin{remark}[一般化について]
    %     $n$-単体（$n\ge3$）では、頂点 $p(i)$ を通り対向面またはその延長と交わる、直線に限らない図形一般をチェビアンと呼ぶ流儀がある。
    %     この時、$n$-単体のチェビアンを特に \emph{$n$-チェビアン}と呼ぶ。
    %     またこの場合対向面またはその延長との交わりが一点とは限らないことに留意してほしい。
    % \end{remark}

    \section{チェバの定理のアフィン性}

    \begin{body}
        \indent 上記の概念を用いて古典的チェバの定理を導出関係を明示しながらアフィン空間の言葉で書き下す。これが正確に古典的チェバの定理の言い換えになっているかどうかは各自確かめて欲しい。
    \end{body}

    \begin{theorem}[古典的チェバの定理の言い換え]

        \noindent\textbf{前提：}\par
        \begin{theoremenumerate}[label=(P\arabic*), leftmargin=4em, topsep=2pt, partopsep=0pt, itemsep=0pt, parsep=0pt]
            \item $k$ は可換体である。
            \item $V$ は $k$-ベクトル空間である。
            \item $P$ は $V$ 上のアフィン空間である。
            \item $\iota = \{0,1,2\}$ である。
            \item 点族$p:\iota \to P$ はアフィン独立で、$A:=p(0), B:=p(1), C:=p(2)$と表す。
            \item $s = \iota$          
            \item 各 $i\in s$ について、有限集合 $F_i:=\iota\setminus\{i\}\subset \iota$ と重み $w_i:\iota\to k$が与えられ、
            \[
            \sum_{j\in F_i} w_i(j)=1
            \]
            が成立し、$D:=(1-t)B+tC, E:=(1-u)C+uA, F:=(1-v)A+vB$と表す。
        \end{theoremenumerate}
        
        \noindent\textbf{仮定：}\par
        \begin{theoremenumerate}[label=(A\arabic*), leftmargin=4em, topsep=2pt, partopsep=0pt, itemsep=0pt, parsep=0pt]
          \item ある点 $p'$ が存在し、$p'$は直線$AD, BE, CF$上に存在する。
        \end{theoremenumerate}
        
        \noindent\textbf{結論：}\par
        \begin{theoremenumerate}[label=(\roman*), leftmargin=4em, topsep=2pt, partopsep=0pt, itemsep=0pt, parsep=0pt]
          \item 以下の等式が成立する：
          \[
          \left(\frac{t}{1-t}\right)\left(\frac{u}{1-u}\right)\left(\frac{v}{1-v}\right)=1
          \]
        \end{theoremenumerate}
    \end{theorem}

    \begin{body}

        \indent 上記の定理は全てアフィン空間の概念で記述されており、線分の長さをはじめとしたユークリッド幾何の概念は一切登場しない。辺の長さの比に見えていたものは重心座標で現れる係数の比であり、
        \[
        \dfrac{BD}{DC} := \dfrac{t}{1-t},\qquad
        \dfrac{CE}{EA} := \dfrac{u}{1-u},\qquad
        \dfrac{AF}{FB} := \dfrac{v}{1-v}
        \]
        である。

        \indent この結果は本稿で作成するモジュールをmathlib4のどのディレクトリのファイルに追加するか決定する際に重要となる。一般化を行うときもユークリッド幾何の計量の概念は避けることが必要である。

    \end{body}

    % =========================================================
    \chapter{チェバの定理の一般化と選定}
    % =========================================================
    \begin{body}
        \indent 本章では上記の議論を踏まえ、一般化チェバの定理の主張として求められる何が相応しいかを明確にする。
    \end{body}

    \section{一般化の方針}

    \begin{body}
        \indent 古典的チェバに対して上記の定式化を行うと定理が様々な観点から一般化できることが見える。

        \indent 例えば以下の観点が考えられる。

        \begin{bodyitemize}
            \item 周辺空間のみ一般化する
            \item $2$-単体を$n$-単体に一般化する
            \item $1$-チェビアンを$2$-チェビアンで考えてみる
            \item チェビアンの数は頂点数より多くとも構わない
        \end{bodyitemize}

        \vskip\baselineskip

        \indent 実際に一般化したものとして以下のようなものが挙げられる。

        \begin{bodydescription}
            \item[\hspace*{1.5em}・候補A：周囲空間のみ一般化]
            $n$次元アフィン空間に三角形を埋め込み、その各頂点から対向する$1$-面上の点へ下ろす$1$-チェビアンを考える．
            これらの直線が（$n$次元アフィン空間内で）共点（または平行）であることと、古典的チェバと同型の比の積等式が成り立つことの必要十分性を主張する定理．
            （本質は古典的チェバの定理と同じで、周囲空間の次元は主張に影響しない．）
            
            \item[\hspace*{1.5em}・候補B：$n$-単体版, 1-チェビアン]
            $n$-単体の各頂点から対向する$(n-1)$-面上の点へ下ろす$1$-チェビアン（直線）を考える．
            これらが共点または平行であることと，同次質量分布型の条件（頂点への質量割当から各対向面上の点が重心として与えられる等）が成り立つことの必要十分性を主張する定理．
            
            \item[\hspace*{1.5em}・候補C：$n$-単体版, 2-チェビアン]
            $n$-単体の各頂点から対向する$(n-1)$-面に対し，$2$次元的データ（面積比・外積・行列式など）で自然に定まる$2$-チェビアンを考える．
            これらが共点または平行であることと，面積比（あるいは行列式）により与えられる条件が成り立つことの必要十分性を主張する定理．
            
            \item[\hspace*{1.5em}・候補D：$n$単体版, multipede]\cite{multipede}
            $n$-単体の各頂点から対向する$k$-面（$1 \leq k \leq n-1$）への$1$-チェビアンを考える（$k$を固定しても，全ての$k$を同時に扱ってもよい）．
            これらが所定の意味で共点（または平行）であることと，対応する係数整合（バリセントリック係数の整合性，または質量分布条件の一般化）が成り立つことの必要十分性を主張する定理．
            一つの頂点から複数のチェビアンが伸びることから$multipede$版チェバと呼ばれる。
                
        \end{bodydescription}

        \indent このような、さまざまな版で一般化されたチェバの定理の総称を\emph{チェバ型定理}と呼ぶ。
    \end{body}

    \section{先行プロジェクトの紹介}
    \begin{body}
        \indent 過去のいくつかのチェバ型定理形式化に関する先行研究を紹介する。

        \begin{bodydescription}
            \item[\hspace*{1.5em}・Lean3時代の形式化]
            期間: $2021$年$12$月 〜 $2023$年$7$月\\
            GitHub: \url{https://github.com/leanprover-community/mathlib3/pull/10632} \cite{mathlib3_ceva}

            Mantas Bakšys氏による試み。
            $2$-単体のチェバの定理のみを形式化、ユークリッド幾何の概念である距離を全面に出した証明を行ったが、mathlib4に追加するには一般性が足りないとして採用されなかった。

            \item[\hspace*{1.5em}・Aristotle.AI（Harmonic社）を用いた形式化]
            期間: $2025$年$12$月 〜 $2026$年$1$月\\
            GitHub: \url{https://github.com/leanprover-community/mathlib4/pull/33388} \cite{mathlib4_ceva_alexeev}

            Boris Alexeev氏による試み。
            Harmonic社が提供するLeanプログラム特化のAIサービスAristotle.AIを用いて形式化した。
            全てのコードがLeanの機械検証を通過し、正確性が保証された。AIの性能向上を示したと言える。
            しかし、$2$-単体のチェバの定理のみを形式化しており、またコミュニティ内の過去の議論を踏まえられていなかったため、mathlib4に追加するには一般性が足りないとして採用されなかった。

            \item[\hspace*{1.5em}・$n$-単体で汎用一般化された初の形式化]
            期間: $2026$年$1$月 〜 $2026$年$1$月\\
            GitHub: \url{https://github.com/leanprover-community/mathlib4/pull/33409} \cite{mathlib4_ceva_myers}

            Joseph Myers氏による試み。
            過去の不十分な事例を踏まえて、$n$-単体に一般化されたチェバの定理順方向を形式化した。
            おそらく現時点で最も一般性が高い形式化である。
            後の章で詳細に見ていく。
            2026/1/14 PR受理。
        \end{bodydescription}

    \end{body}

    \section{本稿で扱う一般化の選定}
    \begin{body}
        \indent 本稿では既存のmathlib4に用意された一般化チェバの定理の順方向の方針を採用する。

        \indent 上述したJoseph Myers氏によって形式化されたものであり、「各頂点と、その対向面または延長上の両者を結ぶ$1$-チェビアンが共点する必要条件」を与える。
    \end{body}

    \subsection{先行研究の主張}

    \begin{body}
        \indent mathlib4内のCeva.leanは主に以下の定理が証明されている。
        
        \begin{bodyenumerate}
            \item (2)を示すための補題
            \item 任意添字点族チェバの定理
            \item 有限添字点族チェバの定理
            \item $2$-単体についての古典的チェバの定理（積等式形）
            \item $2$-単体についての古典的チェバの定理（商等式形）
        \end{bodyenumerate}
        \indent 本稿のテーマに近いのは(3)である。

        \indent (1)で一般化チェバの定理の主張のほとんどを補題として提示、(2)で任意添字点族に対する一般化チェバの定理を示し、その特殊な版として(3)を導いている。
            
        \indent (4),(5)ではさらに(3)の特殊形として3つの添字に対する古典的チェバの定理を積等式と商等式で版を分けて紹介している。

        \indent (3)の主張を敢えて要約すると以下のようになる。
    \end{body}

    \begin{theorem}[チェバ型定理順方向]
        $n$-単体の頂点を$p(0)$,…,$p(n)$とする。
        いくつかの頂点$p(i)$に対し、$p(i)$を含まない任意の面上に点$q(i)$を固定し、直線$p(i)q(i)$を考える。
        こうして得られる複数の直線がある点$p'$で交わる時、$p'$は$p(F_i)\cup \{p(i)\}$が張るアフィン部分空間中に存在して、各iについて$q(i)$の$F_i$が定める点族に関する重心座標、および$p'$の$\iota$が定める点族に関する重心座標は第$i$成分を除いて比例する。
    \end{theorem}

    \begin{body}
        \indent もはや古典的チェバの定理の原形を留めていない。

        \indent 点$q(i)$、直線$p(i)q(i)$それぞれは古典的チェバの定理においてチェビアンの足だったもの、チェビアンだったものである。要するに、「共点が存在すれば、共点とチェビアンの足の重心座標成分の大部分が比例する」ことを主張している。

        \indent 一見古典的チェバとの関連は見えづらい。しかし、以下の点で一般化になっている。

        \indent 一つ目に、\emph{$n$-単体についての主張になっている}。前提の弱化である。$n=2$とすれば三角形についての定理になる。

        \indent 二つ目に、\emph{チェビアンは全ての頂点から下ろす必要がない}。前提の弱化である。直線$p(i)q(i)$がチェビアンに相当する。古典的チェバの定理では、全ての頂点からチェビアンを下ろしている。

        \indent 三つ目に、\emph{チェビアンの足は対向面上でなくて良い}。前提の弱化である。点$q(i)$がチェビアンの足に相当する。古典的チェバの定理では、チェビアンの足は対向面（対辺）上にあった。しかし、この定理中ではチェビアンの足はいずれかの面上にあれば良い。

        \indent 四つ目に、\emph{商の等式条件が係数の比例条件になっている}。結論の強化である。これによりスカラーを体の元に限る必要がなくなる。

        \indent この主張を導出関係を明確にして書き下す。

    \end{body}

    \begin{theorem}[チェバ型定理順方向の言い換え]
        \noindent\textbf{\\前提：}\par
        \begin{theoremenumerate}[label=(P\arabic*), leftmargin=4em, topsep=2pt, partopsep=0pt, itemsep=0pt, parsep=0pt]
            \item $k$ は 環である。
            \item $V$ は $k$-加群である。
            \item $P$ は $V$ 上のアフィン空間である。
            \item $\iota$ は有限の添字集合である。
            \item 点族 $p : \iota \to P$はアフィン独立。
            \item $s \subseteq \iota$ は空でない。
            \item 各 $i\in s$ について，有限集合 $F_i\subseteq \iota$ と重み $w_i:F_i\to k$ が与えられ，
            \[
            i\notin F_i \land \sum_{j\in F_i} w_i(j)=1
            \]
            が成り立つ。
        \end{theoremenumerate}
        \noindent\textbf{\\仮定：}\par
        \begin{theoremenumerate}[label=(A\arabic*), leftmargin=4em, topsep=2pt, partopsep=0pt, itemsep=0pt, parsep=0pt]
            \item ある点 $p'\in P$ が存在し、各 $i \in s$ について、$p'$ は2点 $p(i)$ と $\sum_{j \in F_i} w_i(j)p(j)$ を結ぶ直線上に存在する。
        \end{theoremenumerate}

        \noindent\textbf{\\結論：}\par
        \begin{theoremenumerate}[label=(C\arabic*), leftmargin=4em, topsep=2pt, partopsep=0pt, itemsep=0pt, parsep=0pt]
            \item ある重み $w' : \iota \to k$が存在して
            \[
            \sum_{j \in \iota} w'(j) = 1
            \]
            を満たし、
            \[
            p' = \sum_{j \in \iota} w'(j)p(j)
            \]
            と書ける。
            \item 各 $i \in s$ ごとにスカラー $r_i \in k$ が存在し、次が成り立つ： 
            \[
            \forall i \in s,\ \exists r_i \in k,\ \forall j \in F_i,\ \bigl(w'(j) = r_i\, w_i(j)\bigr).
            \]
        \end{theoremenumerate}
    \end{theorem}

    \begin{proof}  
        (A1) より、任意の $i\in s$ に対して $p'$ は2点 $p(i)$ と $\sum_{j \in F_i} w_i(j)p(j)$ を結ぶ直線上に存在する。
        よって、
        \[
        \exists \alpha_i\in k, p'=\alpha_i\,p(i)+(1-\alpha_i)\sum_{j\in F_i} w_i(j)\,p(j)
        \]
        ここで各 $i\in s$ について、関数 $\bar w_i:\iota\to k$ を次で定める：
        \[
        \bar w_i(j):=
        \begin{cases}
        \alpha_i & (j=i),\\
        (1-\alpha_i)\,w_i(j) & (j\in F_i),\\
        0 & (j\in\iota\setminus(F_i\cup\{i\})).
        \end{cases}
        \]
        （$w_i: F_i\to k$ であり、$\iota$全域で定義されていないため、$j\notin F_i$ の場合を上の定義により$\bar w_i(j)=0$ として埋めた。）
  
        $i\notin F_i$ と $\sum_{j\in F_i}w_i(j)=1$ より
        \[
        \begin{aligned}
        \sum_{j\in\iota} \bar w_i(j)
            &= \bar w_i(i)+\sum_{j\in F_i} \bar w_i(j)+\sum_{j\in\iota\setminus(F_i\cup\{i\})} \bar w_i(j)\\
            &= \alpha_i + (1-\alpha_i)\sum_{j\in F_i} w_i(j) + 0\\
            &= \alpha_i + (1-\alpha_i)\\
            &= 1
        \end{aligned}
        \]
        が成り立つ。また上の定義から
        \[
        p'=\sum_{j\in\iota} \bar w_i(j)\,p(j)
        \]
  
        以上より、各 $i\in s$ に対して
        \[
        \sum_{j\in\iota} \bar w_i(j)=1,\qquad
        p'=\sum_{j\in\iota} \bar w_i(j)\,p(j)
        \]
        が得られた。
  
        また(P5)より、$p:\iota\to P$はアフィン独立であるから$\AffHull(p)$中の点を点族$p$のアフィン結合で一意に表せる。
        (A1)より、$p'\in P$である。
        従って任意の $i,i'\in s$ について
        \[
        \bar w_i = \bar w_{i'}
        \]
  
        よって、任意に $i_0\in s$ を一つ固定して
        \[
        w' := \bar w_{i_0}:\iota\to k
        \]
        とおけば、全ての $i\in s$ に対して $w'=\bar w_i$。
  
        したがって
        \[
        \sum_{j\in\iota} w'(j)=1,\qquad
        p'=\sum_{j\in\iota} w'(j)\,p(j)
        \]
        (C1) が示された。
  
        また、$r_i:=1-\alpha_i$ とおけば、定義より$w'=\bar w_i$ なので、任意の $j\in F_i$ について
        \[
        \begin{aligned}
            \forall j\in F_i, 
            w'(j) &=\bar w_i(j)
                &=(1-\alpha_i)\,w_i(j)
                &=r_i\,w_i(j)
        \end{aligned}
        \]
        (C2)が示された。
    \end{proof}

    \begin{body}
        \indent これが古典的チェバの定理の一般化になっていることは、前提・仮定・結論に含まれる命題の強弱を比較すれば直ちにわかる。

        \indent mathlib4内のCeva.Leanでは上記の主張をベースにLean4向けに整形した証明がなされている。

        \indent この定理を参考に、本稿では古典的チェバの定理の逆を一般化して示し形式化する。
    \end{body}

    % =========================================================
    \chapter{Lean4での実装}
    % =========================================================

    \section{チェバ型定理逆方向の主張と自然言語証明}
    \begin{body}
        古典的チェバの定理の逆の主張は導出関係を明確にすれば以下のような形になる。
    \end{body}

    \begin{theorem}[古典的チェバの定理逆の言い換え]
        \noindent\textbf{\\前提：}\par
        \begin{theoremenumerate}[label=(P\arabic*), leftmargin=4em, topsep=2pt, partopsep=0pt, itemsep=0pt, parsep=0pt]
            \item $k$ は可換体である。
            \item $V$ は $k$-ベクトル空間である。
            \item $P$ は $V$ 上のアフィン空間である。
            \item $\iota = \{0,1,2\}$ である。
            \item 点族$p:\iota \to P$ はアフィン独立で、$A:=p(0), B:=p(1), C:=p(2)$と表す。
            \item $s = \iota$          
            \item 各 $i\in s$ について、有限集合 $F_i:=\iota\setminus\{i\}\subset \iota$ と重み $w_i:\iota\to k$が与えられ、
            \[
            \sum_{j\in F_i} w_i(j)=1
            \]
            が成立し、$D:=(1-t)B+tC, E:=(1-u)C+uA, F:=(1-v)A+vB$と表す。
        \end{theoremenumerate}
      
        \noindent\textbf{仮定：}\par
        \begin{theoremenumerate}[label=(A\arabic*), leftmargin=4em, topsep=2pt, partopsep=0pt, itemsep=0pt, parsep=0pt]
            \item 係数 $t,u,v$ が
            \[
            \left(\frac{t}{1-t}\right)\left(\frac{u}{1-u}\right)\left(\frac{v}{1-v}\right)=1
            \]
            を満たす。
        \end{theoremenumerate}
      
        \noindent\textbf{結論：}\par
        \begin{theoremenumerate}[label=(\roman*), leftmargin=4em, topsep=2pt, partopsep=0pt, itemsep=0pt, parsep=0pt]
            \item 点$p'$が存在して
        \end{theoremenumerate}
    \end{theorem}

    \begin{body}
        \indent この主張を元に、順方向同様一般化を行う。
    \end{body}
    

    \begin{theorem}[チェバ型定理逆方向]
        \noindent\textbf{\\前提：}\par
        \begin{theoremenumerate}[label=(P\arabic*), leftmargin=4em, topsep=2pt, partopsep=0pt, itemsep=0pt, parsep=0pt]
            \item $k$ は体である。
            \item $V$ は $k$-ベクトル空間である。
            \item $P$ は $V$ 上のアフィン空間である。
            \item $\iota$ は有限の添字集合である。
            \item $s \subseteq \iota$ は空でない。
            \item 点族 $p:\iota\to P$はアフィン独立。
            \item 各 $i\in s$ について，重み $w_i:\iota\to k$ が与えられ，次を満たす：
            \[
            \sum_{j\in\iota} w_i(j)=1,\qquad w_i(i)=0,\qquad q(i)=\sum_{j\in\iota} w_i(j)\,p(j).
            \]
        \end{theoremenumerate}
    
        \noindent\textbf{\\仮定：}\par
        \begin{theoremenumerate}[label=(A\arabic*), leftmargin=4em, topsep=2pt, partopsep=0pt, itemsep=0pt, parsep=0pt]
            \item ある写像$w:\iota\to k$ が存在して、次の条件を満たす：
            \[
            \forall i\in s,\ \exists r_i\in k,\ \forall j\in\iota,\ \bigl(j\neq i \Rightarrow w(j)=r_i\,w_i(j)\bigr).
            \]
            \item $w$ は零写像ではない。
            \[
            \exists j\in\iota,\ w(j)\neq 0.
            \]
        \end{theoremenumerate}
  
        \noindent\textbf{\\結論：}\par
        \begin{theoremenumerate}[label=(C\arabic*), leftmargin=4em, topsep=2pt, partopsep=0pt, itemsep=0pt, parsep=0pt]
            \item $\sum_{j\in\iota} w(j)\neq 0$ のとき、$p'\in P$が存在して、任意の $i\in s$ について $p'$ は2点 $p(i)$ と $q(i)$ を結ぶ直線上に存在する。
            % \[
            % w^{\mathrm{fin}}(j):=\frac{w(j)}{\sum_{t\in\iota} w(t)}\quad (j\in\iota)
            % \]
            % とおけば $\sum_{j\in\iota} w^{\mathrm{fin}}(j)=1$ が成り立つ。さらに
            % \[
            % p':=\sum_{j\in\iota} w^{\mathrm{fin}}(j)\,p(j)
            % \]
            \item $\sum_{j\in\iota} w(j)=0$ のとき、任意の$i\in s$について直線 $p(i)q(i)$は互いに平行である。
            % \[
            % \forall i_0\in \iota, v^{\infty}:=\sum_{j\in\iota} w(j)\,\bigl(p(j)\dotminus p(i_0)\bigr)\in V
            % \]
            % とおくと，$v^{\infty}\neq 0$ が成り立つ。さらに任意の $i\in s$ についてある $t_i\in k$ が存在して
            % \[
            % q(i)=(t_i\cdot v^{\infty})\dotplus p(i)
            % \]
            % が成り立つ。
        \end{theoremenumerate}
    \end{theorem}

    \begin{proof}
        任意の$i\in s$をとる。
        $\sum_{j\in\iota} w(j)\neq 0$より、関数 $w':\iota\to k$ を次で定める：
        \[
        w'(j):= \dfrac{w(j)}{\sum_{t\in \iota} w(t)}
        \]
        定義より明らかに、$\sum_{t\in \iota} w(t) = 1$
        この重みを用いて点族$p$のアフィン結合を考える。
        \[
        p'=\sum_{j\in \iota} w'(j)p(j)
        \]
    \end{proof}

    \begin{body}
        \indent 上記の定理は要するに、「
    \end{body}


    \section{設計}
    \subsection{設計の問題点}
    \begin{body}
        \indent 素朴には上記の定理をそのままLean4に書き直せば良いと考えられる。しかしここで生じる問題は、$3$章で行った概念の定義は本稿で証明を行いやすくするために採用したものであり、mathlib4中に用意された概念の定義とはいささか異なる。例えば、チェバ型定理という単一の定理を紙面上で示すと言う目的でアフィン結合$\to$アフィン部分空間という定義の流れを本稿では選択したが、これは$V$が$k$-加群であることを先に仮定しているという意味で数学的にあまり綺麗な方法ではない。実際mathlib4中の定義ではアフィン部分空間$\to$アフィン結合の流れが採用されており、本稿で用いた定義とは流れも見た目も異なっている概念も多い。よって、形式化の際はこのような定義の読み替えを適宜行いつつ翻訳を行っていく必要がある。
    \end{body}

    \subsection{設計方針と採用理由}
    \begin{body}
        \indent 前章で見たようにJoseph Myers氏は最も一般的な定理から徐々に特殊な場合の定理を導くという流れをとっており、チェバ型定理逆方向の実装でも同様の方法を取ることにする。

        \indent 具体的には、以下の通りである。

        \begin{bodyenumerate}
            \item (2)を示すための補題
            \item 任意添字点族チェバの定理の逆
            \item 有限添字点族チェバの定理の逆
            \item 有限添字点族チェバの定理の逆（平行を除いた版）
            \item $2$-単体についての古典的チェバの定理の逆
            \item $2$-単体についての古典的チェバの定理の逆（平行を除いた版）
        \end{bodyenumerate}

        % \indent 基本的に構成は同じであり、

    \end{body}

    \section{追加モジュールの実装}
    \subsection{概念と既存との対応}
    \begin{body}
        \indent $3$章で紹介した数学的概念、補題がmathlib4上でどのように実装されているのか紹介する。
        \begin{bodydescription}[style=nextline,leftmargin=4.5em,labelsep=1em,font=\normalfont\bfseries,itemsep=0pt, topsep=2pt, parsep=0pt]
          \item[アフィン空間] \nolinkurl{AddTorsor V P}, \nolinkurl{AffineSpace V P}
          \item[アフィン独立] \nolinkurl{AffineIndependent k p}
          \item[アフィン結合] \nolinkurl{Finset.affineCombination k s p w}
          \item[重心座標] \nolinkurl{AffineBasis.coord}
          \item[アフィン包] \nolinkurl{affineSpan k (s : Set P)}
          \item[凸包] % TODO
          \item[単体] \nolinkurl{Affine.Simplex k P n}
          \item[面] \nolinkurl{Affine.Simplex.face}
          \item[対向面] \nolinkurl{Affine.Simplex.faceOpposite}
        \end{bodydescription}

        \indent 上記のモジュールを使いながら、不足しているものを適宜追加し実装を進める。
    \end{body}

    \subsection{追加するモジュール}
    \begin{body}
        
    \end{body}

    % \subsection{実装の工夫事項・特記事項}
    % \begin{bodydescription}
    %     \item \textbf{実装を急いでいます}
    % \end{bodydescription}

    % =========================================================
    \chapter{結論と今後の展望}
    % =========================================================

    \section{結論}
    \begin{body}
        \indent 既にJoseph Myers氏によって形式化された一般化チェバの定理を参考に、一般化チェバの定理の逆を形式化し、mathlib4にPRを提出した。
    \end{body}

    \section{今後の課題}
    \begin{body}
        \indent 当初は一般化チェバの定理の順方向を実装する予定だったが、コミュニケーション不足によりJoseph Myers氏に先行されることになった。研究内容の共有がいかに大切かを痛感した経験となった。

        \indent また、数学にしろ情報学にしろ、何を達成したいかを正確に把握して最適な言語・理論を選択することが不可欠であることを痛感した。単一の定理の証明を紙面上で構築するための概念の定義なのか理論の一部としてライブラリに用意する概念の定義なのかという違い、自然言語で定義するのか一階述語論理で定義するのか型理論で定義するのかという違い、によって最適な定義が揺れるためその選定に最も苦戦した。
        
        \indent 今回培ったOSS開発の作法に従い、次回以降のプロジェクトにも取り組んでいきたい。
    \end{body}

    \section{Lean4の普及と形式化研究の展望}
    \begin{body}
        \indent 現在LeanのAutoformalizationという技術開発が進んでおり、自然言語で記述した数学の証明をそのままLeanの形式言語に書き換えることが目指されている。\\
        今回は個別の定理を形式化すると言うプロジェクトであったが、この技術を用いればある数学的な証明を自然言語で記述するだけで、立ちどころに形式証明が得られると考えられている。
        この技術は、Lean研究を飛躍的に促進すると考えられる。個別の定理の形式化だけでなく、このAutoformalizationに関する見識も深めていきたい。
    \end{body}

    % =========================================================
    \appendix
    % =========================================================

    % =========================================================
    % 付録：型理論についての補足
    % =========================================================
    \chapter{型理論についての補足}\label{appendix:type-theory}

    \section{本付録の目的}
    \begin{body}
        \indent 本章で触れたように型理論では、命題を型・証明をその型の項、として表現する。しかしこれだけでは主に一階述語論理ベースの集合論に慣れている読者にはその全貌は見せられないように思う。よってよりイメージのつきやすい型理論の説明を一階述語論理と比較し、形式体系とは何かに立ち返りながら行うことが本付録の目的である。
    \end{body}

    \section{形式体系とは}

    \begin{body}
        \indent 数学の推論を厳密に形式化して分析する分野を数理論理学と呼び、形式体系とはその中で用いる概念の一つである。数理論理学は構文論と意味論に分かれるが、両者は形式体系を別のアプローチで分析する。構文論とは記号列の操作（導出）の規則に着目する手法であり、意味論とは記号列に解釈を与えてその真偽を論じる手法である。

        \indent 一階述語論理や高階述語論理、型理論など、形式体系と呼ばれるものは様々あり、形式言語を用いて数学の証明・推論を様々な方法で表現する。構文論的な観点で見たときこれらの多くに共通する眼目は、証明という対象を有限のデータとして扱う点にある、つまり\emph{証明とは有限の記号列でなければならない}。
    
        \indent 形式体系はしばしば \emph{言語}、\emph{公理}、\emph{推論規則}の三つ組として記述される。言語は「許される記号と、それらから整式（文法を満たす記号列）を作る規則」を与える。公理は「証明の出発点として採用する式またはその集合」である。推論規則は「前提となる式から結論となる式を得る規則」である。これら自体は有限である必要はない。実際、変数記号は自然数の個数だけ無数に用意できるし、公理はスキームという形で実質無限本存在する、その上推論規則の中には$\omega$規則$\frac{\{\varphi(n)\}_{n\in\mathbb{N}}}{\forall n\in\mathbb{N}\,\varphi(n)}\;(\omega)$と呼ばれるものもあり、推論が無限を含んでいる。

        \indent しかし証明を構築するときは話が別である。証明とは「人間が数学的主張の正しさを検証するための道具」であるから有限でなくては検証ができない。よっていざ証明を構築するときは、言語から有限個の記号を選び、公理スキームの有限個を利用し、有限個の仮定から有限個の結論が得られる推論規則のみを利用する。したがって形式体系を構築するにあたって、言語の無限個の記号、公理の無限本の公理図式は許容されるが、$\omega$規則のような\emph{無限推論規則}は「証明が有限」という前提自体を捨てる必要が生まれるため一般的な形式体系の議論では許容されない。（これを認めた無限論理・無限証明論と呼ばれる分野も存在するが、無限を扱える代わりに完全性やコンパクト性などの豊かな性質の多くが失われるまたは変質してしまう別物として区別され研究されている）
    
        \indent 本稿では、与えられた公理（集合）に推論規則を適用して得られる導出可能な式全体の閉包を \emph{理論体系}と呼ぶ。「これらの式が何を意味するか」という解釈を与えるのは意味論に譲る。
    
        \indent 例えば、多くの数学書が暗黙のうちに議論に用いている集合論（典型的には ZF / ZFC）は、形式体系としては「集合 $\in$ を含む言語」「公理（外延性・分出・置換スキーム等）」「古典論理の推論規則」を備えた理論として以下のように整理できる。
        
        \vskip\baselineskip

        \noindent\textbf{集合論}
        \begin{bodyitemize}[leftmargin=7em]
            \item[言\hspace*{2em}語] : 一階述語論理記号$\cup$変数記号$\cup$補助記号$\cup\in$と論理式の定義
            \item[推論規則] : 自然演繹
            \item[公\hspace*{2em}理] : ZF / ZFC
        \end{bodyitemize}

        \vskip\baselineskip

        \indent まず、言語を定義する。記号集合$\{\neg,\to,\forall,\exists$,…$\}\cup \{v_{0},v_{1},v_{2}$,…$\}\cup \{\in\}$を用意し、この集合から$\neg v_{2}\forall, \forall v_{0}(v_{0}=v_{0}), \cup v_{n}\cap v_{n-1}=v_{0}$のように無数の記号列を構成できる。そして得られた記号列から記号列として認めるものを定義する。認めた記号列を論理式と呼び$\phi,\psi,\varphi$などと略記する。先ほどの例では$\forall x(x=x)$のみ論理式である。\\
        次に、推論規則を定義する。有名な例では含意除去$\infer[\to E]{\varphi \\ \varphi \to \psi}{\psi}$（Modus Ponensのこと）などがある。先ほどの例に対しては例えば全称除去$\infer[\forall E]{\forall x\,\varphi(x)}{\varphi(t)}$を用いて$v_{0}=v_{0}$などの新たな論理式が得られる。\\
        最後に公理を決める。Kunenの定義に従えば、外延性の公理
        \[
        \forall v_0\,\forall v_1\;\Bigl(\forall v_2\,(v_2\in v_0 \leftrightarrow v_2\in v_1)\;\rightarrow\; v_0=v_1\Bigr)
        \]
        などがある。論理式の中で証明の出発点として採用する特別な論理式である。

        \indent こうして得られた公理から推論規則による変形を通して得られたものだけが定理と呼ばれ、理論体系を構成する一員となる。

        \indent こう書くと、自然に「では言語、公理、推論規則のいずれかを変更したらどうなるだろう？」という疑問が自然に湧いてくるであろう。公理的集合論では主にこの公理の部分を様々に入れ替え、付け加え、除くことで理論体系がどう変化するかを分析する。また、推論規則としてはヒルベルト式やシーケント計算が存在し定理の証明のしやすさに影響を与える。

        \indent 同様にしてこの三つ組を変更して得られるのが、型理論である。（正確にいうと、上記の三つ組は一階述語論理ベースの形式体系を分類するには有要だが、型理論ベースの形式体系を分類するには不向きである。その上で敢えて同じ視点でCICを眺めてみる。）

        \textbf{ここからほぼAIです}
        \vskip\baselineskip
        \begin{bodyitemize}[leftmargin=4em]
            \item[\textbf{CIC}] \mbox{}
            \item[言語]
                文脈$\Gamma$・項$t$・型$A$からなる依存型付き$\lambda$計算の構文
            \item[公理]
                特になし
                ただし必要に応じて定数宣言$c:A$を追加して公理的拡張を行うことがある。
            \item[推論規則]
                判断
                \[
                \Gamma \ \mathsf{ctx},\qquad
                \Gamma \vdash A \ \mathsf{type},\qquad
                \Gamma \vdash t : A
                \]
                を生成する型付け規則（変数規則・$\Pi$の形成/導入/除去・帰納型の形成/導入/除去など）、計算規則（$\beta$簡約・帰納型の$\iota$計算規則等）および、変換規則
        \end{bodyitemize}

        \indent 見慣れない用語が並びイメージが付きづらいように思う。一階述語論理に慣れている読者が型理論をゼロから捉えるうえで重要なのは、集合論のように「ある宇宙をまず仮定して、その上で式の真偽を語る」という発想をいったん脇に置くことである。型理論（少なくともCICの系譜）では、最初にあるのは\emph{対象世界}ではなく、\emph{判断を生成する推論体系}である。つまり「何が式（項・型）として許されるか」「その式がどの型を持つと主張してよいか」「二つの式が同一視されてよいか」を、構文と規則によって定義する。集合論が「$\in$ を含む一階言語 + 公理 + 推論規則」で理論を作るのと同様に、型理論は「判断（judgement）+ 生成規則」で理論を作る。

        \indent まず\emph{判断}の形だけを固定する。判断とはメタ主張のことを指す。CICの核は典型的に次の三種の判断で動く：
        (1) 文脈が整っている $\Gamma\ \mathsf{ctx}$、
        (2) 文脈の下で型が整っている $\Gamma\vdash A\ \mathsf{type}$、
        (3) 文脈の下で項が型を持つ $\Gamma\vdash t:A$。
        ここで「文脈」は自由変数とその型の並び（および場合によっては定義）であり、一階論理の「変数割当」や「仮定集合」に対応する。だが対応は完全ではない。型理論では文脈そのものが推論の対象であり、文脈拡張（$\Gamma,x:A$）が推論規則として明示される。

        \indent 次に、\emph{構文の生成}を与える。依存型付き$\lambda$計算では、項は変数、適用 $t\,u$、抽象 $\lambda x:A.\,t$ 等から作られ、型は関数型（依存積）$\Pi x:A.\,B$ 等から作られる。ここで本質は「型が項に依存できる」ことで、$\Pi x:A.\,B$ に現れる $B$ は $x$ を含みうる。これが、一階論理で言えば「述語が変数に依存する」ことに対応するが、型理論では依存が\emph{型形成のレベル}で起きる点が決定的に違う。

        \indent しかし構文だけでは何も主張できない。そこで三段階目として、\emph{型付け規則（判断生成規則）}を置く。最も基本的なのは変数規則で、$\Gamma$ が整っていて $\Gamma$ に $x:A$ が含まれるなら $\Gamma\vdash x:A$ を得る。続いて、$\Pi$型について形成・導入・除去（いわゆる関数型の形成、ラムダ抽象、適用）が与えられる。これらは自然演繹の対応物として読むのが早い：$\Pi$は全称量化と含意を統合したような振る舞いをし、$\lambda$は仮定導入、適用は仮定消去に相当する。ただし「論理式の証明」という読みは\emph{後から可能になる解釈}であって、体系の定義そのものはあくまで型付け規則の集合として完結する。

        \indent 四段階目が、集合論側の読者が最初に転ぶ地点である。型理論には「等号」が二種類ある。ひとつは\emph{定義的等しさ}（definitional equality）で、計算により同一視してよいという約束である。$\beta$簡約（$(\lambda x.\,t)\,u \leadsto t[x:=u]$）や帰納型の計算規則（$\iota$規則）がここに入る。もうひとつは\emph{命題としての等しさ}（identity type）で、$t=u$ を証明（項）として与える対象である。CICの核（カーネル）が自動で判定するのは前者であり、後者は通常の命題と同様に証明が必要になる。ここを混同すると、型理論が「計算の体系」なのか「論理の体系」なのか見失う。実態は両方で、計算が\emph{型付けの同値判定}に組み込まれている点が特徴である。

        \indent 五段階目として、\emph{帰納型}を導入する。CICが「表現力の強い数学の言語」たりうるのは、自然数・リスト・木・有限集合・（適切な形の）集合様構造などを、帰納型の形成規則とその帰納法/再帰原理として内蔵するからである。ここで公理を足すのではなく、\emph{新しい型構成子とその推論規則}を足す、という設計思想が見える。集合論で「分出・置換・無限」等を公理として積むのに対して、型理論では「このデータ型をこういう生成子で作る」「それに対する関数は再帰で定義できる」という規則として積む。

        \indent 六段階目が\emph{宇宙}である。素朴に「型の全体」を一つの型として置くとラッセル型の破綻が起きるため、CICでは $Type_0 : Type_1 : Type_2 : \cdots$ のような階層を置き、型がどの宇宙に属するかを追跡する。これは集合論の累積階層 $V_\alpha$ を思い出せば直観は掴めるが、同一ではない。宇宙は「型の型」というメタ構造を\emph{対象言語内に反映するための装置}であり、同時に矛盾回避のための層構造である。

        \indent ここまでで、公理が「特になし」と書かれている意味が見えてくる。CICは、何らかの特定の数学的対象（例えば $\mathbb{N}$ の標準モデル）を最初に仮定しない。代わりに、判断生成規則が許す限りの構成を行う。その上で、必要なら定数宣言 $c:A$ を追加して拡張する。これは集合論で言えば「新しい定数記号を導入し、追加公理でそれを規定する」操作に似ている。ただしCICでは、その追加は「証明できないが仮定したいもの（例えば選択公理、古典命題の排中律、特定の外延性原理など）」を\emph{公理として型に住まわせる}形で行われる。どの公理を足すかで、計算性・構成性・抽出可能性といった性質が大きく変わるので、無自覚に足すと設計が壊れる。

        \indent 以上を踏まえると、CICは「一階論理 + $\in$」のような単一の基礎記号から世界を立てる流儀ではなく、「判断の生成規則（型付け）+ 計算（簡約）+ 帰納原理」という三点セットから世界を立てる流儀だと言える。集合論が\emph{真理値}を中心に体系を組むのに対し、型理論は\emph{構成（項）}を中心に体系を組む。命題＝型、証明＝項という対応はこの上に乗る見方であり、最初に置くべき土台は「判断生成規則の体系」である。

        \indent 型理論は一階述語論理とは異なる観点から定義される形式体系を提供し、数学の証明とは何かを探る重要なツールとなっている。
    \end{body}

    % =========================================================
    % 付録：　古典的チェバの定理の版の教科書による違いについて
    % =========================================================
    \chapter{古典的チェバの定理の版の教科書による違いについて}\label{appendix:classical-ceva-version}

    \section{本付録の目的}
    \begin{body}
        \indent 本付録では、高校数学で学ぶ古典的チェバの定理の版の違いがなぜ生まれ、どのような意図があるのかを解説する。
    \end{body}

    \section{古典的チェバの定理の版の種類}
    \begin{body}
        \indent 古典的チェバの定理は大きく分けて以下の3つの形で紹介されうる。
        \begin{bodyitemize}
            \item[①] $D, E, F$を辺上に限定した、3直線が1点で交わるための必要十分条件
            \item[②] $D, E, F$を辺の延長上でも許した、3直線が1点で交わる必要条件(\hyperref[thm:classical-ceva]{定理3.1})
            \item[③] $D, E, F$を辺の延長上でも許した、3直線が1点で交わる必要十分条件
        \end{bodyitemize}

        \indent 一見それぞれ異なる主張をしているように見えるがこれらは数学的に何も問題はなく、\emph{背後の幾何学体系に何を採用するか}で主張の範囲や述べ方が違うだけである。

    \end{body}

    \section{古典的チェバに版の違いが生まれる要因}
    \begin{body}
        \indent ①〜③の中で本稿で取り扱っている②のみ必要十分性を主張する定理になっていない。試しに定理の安直な逆命題を考えてみると成立しないことがわかる。その原因となっているのが、下図のような、チェビアンが平行になる場合である。
    \end{body}

    \begin{center}
        \begin{tikzpicture}[scale=0.8, x=1cm,y=1cm, line cap=round, line join=round, >=Latex]
            % ---- keep the whole picture nicely centered ----
            \path[use as bounding box] (-1.2,-1) rectangle (7.2,10.2);
    
            % --- Title (Added per request) ---
            \node[font=\bfseries] at (3,9.8) {(図)等式条件を満たしながらチェビアンが平行となる場合};
    
            % --- coordinates ---
            \coordinate (A) at (0,0);
            \coordinate (B) at (6,0);
            \coordinate (C) at (4,3);
    
            \coordinate (D) at (0,9);
            \coordinate (E) at (6,4.5);
            \coordinate (F) at (4,0);
    
            % --- triangle ---
            \draw[thick] (A) -- (B) -- (C) -- cycle;
    
            % --- side extensions (dashed) ---
            \draw[dashed] (B) -- (D);  % extension of BC to meet x=0
            \draw[dashed] (A) -- (E);  % extension of CA to meet x=6
    
            % --- points ---
            \fill (A) circle (1.6pt) node[below left] {$A$};
            \fill (B) circle (1.6pt) node[below right] {$B$};
            \fill (C) circle (1.6pt) node[above, xshift=2pt] {$C$};
    
            \fill (D) circle (1.6pt) node[left] {$D$};
            \fill (E) circle (1.6pt) node[right] {$E$};
            \fill (F) circle (1.6pt) node[below] {$F$};
    
            % --- Improved Parallel Mark Style ---
            \tikzset{
              parallelArrow/.style={
                postaction={decorate},
                decoration={
                  markings,
                  mark=at position 0.5 with {
                    \draw[line width=0.8pt] (-3pt,-3pt) -- (0pt,0pt) -- (-3pt,3pt);
                  }
                }
              }
            }
    
            % --- cevians (parallel) with fixed arrows ---
            \draw[very thick, parallelArrow] (A) -- (D);
            \draw[very thick, parallelArrow] (B) -- (E);
            \draw[very thick, parallelArrow] (F) -- (C); 
    
        \end{tikzpicture}
    \end{center}

    \begin{body}
        \indent この場合、等式条件は成立するがチェビアンが平行となってしまい共点していない。まさにこの問題によって、背景とする幾何学体系による定理の述べ方が変わるのである。
    \end{body}

    \section{各々の版の背景}

    \begin{body}
        \indent 結論から述べると、①〜③はそれぞれユークリッド幾何、アフィン幾何、射影幾何の観点から古典的チェバの定理を述べたものである。

        \indent ①はユークリッド幾何の観点で述べた形である、$BD$、$DC$などは全てユークリッド距離であり常に正となる。ユークリッド距離ではチェビアンの足が対辺を外分するときをうまく表現できないため、あらかじめチェビアンの足は対辺中に限定される。チェビアンが平行になることは排除できないが、その場合等式条件が成立しないため必要十分性を主張する定理となる。

        \indent ②はアフィン幾何の観点で述べた形である。ここで登場するのはアフィン比と呼ばれる$\dfrac{BD}{DC}$という値であり、辺を成す$2$頂点の重心座標から得られるものである。高校範囲内ではアフィン幾何を登場させられないため、ユークリッド幾何の範疇で定義できる有向線分という考え方を用いて表現している。しかしこの場合定理の逆がそのままでは成り立たず、等式条件が成立してもチェビアンが平行になってしまうという反例が生じる。

        \indent ③は射影幾何の観点で述べた形である。射影幾何はユークリッド幾何やアフィン幾何を包含する幾何学の大きな枠組みの一つである。重要な特徴が空間内に無限遠点を含むため、平行という概念を「無限遠点で交点を持つ」と言える点である。定理の主張自体は②と変わっていないが、交点をもつという言葉の意味が広がったことで述べ方が変わっている。

        \indent 高校数学で習うのは高々ユークリッド幾何であり違いを説明するのが困難なため、追加の前提条件を加えるなどして主張が偽にならないように工夫して解説するという教育的な配慮がなされている。

        \indent 本稿ではアフィン幾何に基づいて議論を進めた。「チェビアンが平行でない」という趣旨の前提を加えれば強引に必要十分性を主張する定理にできるが、定理の自然さが損なわれると考え②の形を選択した。
    \end{body}

    \backmatter
    \addchaptertotoc{謝辞}
    \begin{body}
        % TODO
    \end{body}

    \begin{thebibliography}{99}
        \bibitem{Kunen}Kenneth Kunen, \emph{Set Theory: An Introduction to Independence Proofs}, North-Holland (Elsevier Science Publishers B.V.), 1980.
        \bibitem{LeanProverCommunity} Lean Prover Community, \emph{Lean Prover Community}, \url{https://leanprover-community.github.io/}, accessed 2026-01-18.
        \bibitem{DeepMindIMO2025} Google DeepMind, \emph{Advanced version of Gemini with Deep Think officially achieves gold-medal standard at the International Mathematical Olympiad}, 2025.
        \bibitem{multipede} Dov Samet, \emph{An Extension of Ceva's Theorem ton-Simplices}, 2021.
        \bibitem{LeanResearch} 秋田 隼, \emph{LeanResearch}, GitHub repository, \url{https://github.com/Aj1905/LeanResearch.git}, 2026.
        \bibitem{mathlib3_ceva} Mantas Bakšys, \emph{Ceva's Theorem (mathlib3)}, GitHub Pull Request \#10632, \url{https://github.com/leanprover-community/mathlib3/pull/10632}, 2021--2023.
        \bibitem{mathlib4_ceva_alexeev} Boris Alexeev, \emph{Ceva's Theorem (mathlib4)}, GitHub Pull Request \#33388, \url{https://github.com/leanprover-community/mathlib4/pull/33388}, 2025--2026.
        \bibitem{mathlib4_ceva_myers} Joseph Myers, \emph{Generalized Ceva's Theorem in $n$-dimensional affine space (mathlib4)}, GitHub Pull Request \#33409, \url{https://github.com/leanprover-community/mathlib4/pull/33409}, 2026.

        \vskip\baselineskip

        また、研究にあたりOpenAI社のChat GPT 5.2 thinkingモデルを大いに活用した。以下は使用したチャット履歴の一部である。\\
        \url{https://chatgpt.com/g/g-p-696f0c41aa3c8191928eeed4b3fe9bb8-mathlib4-affinespace/project}
    \end{thebibliography}

\end{document}
